\documentclass[twocolumn]{aastex631}
\begin{document}
\title{The reionization ionizing-photon-budget shortfall for star-forming galaxies is not robust to systematics at z\~{}5}
\author{NebulaMind Lab (autonomous pipeline)}
\affiliation{NebulaMind Lab, \url{https://lab.nebulamind.net}}
\begin{abstract}
Reionization ionizing-photon-budget reconciliation at z\~{}5: star-forming galaxies require f\_esc=0.019 (+0.020/-0.010) to close the budget (Madau-Dickinson SFRD, log xi\_ion=25.5+/-0.15, clumping C=2-5, JWST-SFRD tail), versus indirect-proxy-inferred f\_esc=0.062 (+0.108/-0.039) from LzLCS O32/beta calibrations. Median delta(required-inferred)=-0.039 dex-frac (16-84\%: -0.148 to +0.001); 17\% of the systematic MC shows a shortfall. Conclusion: the budget CLOSES within the systematic, and the sign holds under both O32 and beta calibrations. The result is bounded by the xi\_ion x clumping x proxy-calibration systematic, not by statistics. Generated autonomously from public data (jwst) via the NebulaMind Lab runner: a bounded, reproducible, descriptive study.
\end{abstract}
\section{Introduction}
Recent studies have highlighted a potential crisis in the reionization photon budget, suggesting that star-forming galaxies may not produce enough ionizing photons to account for the observed reionization process [Muñoz2024]. This discrepancy has sparked interest in reconciling the ionizing-photon-budget using literature-anchored calculations. Previous research has explored various aspects of this issue, including the role of galaxy ionizing photon budgets at high redshifts [Duncan2015] and the challenges posed by absorption-dominated reionization scenarios [Davies2021]. However, a systematic reconciliation of these factors is necessary to better understand the reionization process.
\section{Data and method}
To address this, we employed a literature-anchored budget calculation that does not rely on survey catalog data. Instead, we used the cosmic star formation rate density (SFRD) from Madau \& Dickinson's [Madau2017] analytic fitting function, along with published values for xi\_ion and O32/beta f\_esc proxy calibrations from LzLCS studies [Chisholm+22, Flury+22; Simmonds+24]. By adopting these established values, we aimed to systematically reconcile the reionization ionizing-photon-budget using a method that focuses on the ionizing-photon-budget.
\section{Result}
Our calculation reveals that star-forming galaxies require an escape fraction of f\_esc=0.019 (+0.020/-0.010) to close the budget at z\~{}5, given the Madau-Dickinson SFRD, log xi\_ion=25.5±0.15, clumping C=2-5, and JWST-SFRD tail. This value is compared to the indirect-proxy-inferred f\_esc=0.062 (+0.108/-0.039) derived from LzLCS O32/beta calibrations. The median delta between required and inferred values is -0.039 dex-frac (16-84\%: -0.148 to +0.001), with 17\% of the systematic Monte Carlo showing a shortfall.
\begin{figure}\centering\includegraphics[width=\columnwidth]{result.png}\caption{The reionization ionizing-photon-budget shortfall for star-forming galaxies is not robust to systematics at z\~{}5}\end{figure}
\section{Caveats}
It is essential to acknowledge the limitations of our approach. Our analysis relies on an automated, single-selection, uncalibrated measurement, which may not fully capture the complexities of the reionization process. The accuracy of our result depends heavily on the adopted literature values and calibrations, which are subject to their own uncertainties and potential biases. Furthermore, our method does not account for other factors that could influence the ionizing-photon-budget, such as variations in galaxy properties or environmental effects. As a result, while our study provides valuable insights into the reionization photon budget crisis, further research is needed to refine our understanding of this complex phenomenon.
\section*{References}
\footnotesize
[Muoz2024] Muñoz et al. (2024). Reionization after JWST: a photon budget crisis?. bibcode 2024MNRAS.535L..37M \par [Park2022] Park et al. (2022). Calibrating excursion set reionization models to approximately conserve ionizing photons. bibcode 2022MNRAS.517..192P \par [Duncan2015] Duncan et al. (2015). Powering reionization: assessing the galaxy ionizing photon budget at z < 10. bibcode 2015MNRAS.451.2030D \par [Davies2021] Davies et al. (2021). The Predicament of Absorption-dominated Reionization: Increased Demands on Ionizing Sources. bibcode 2021ApJ...918L..35D \par [Madau2017] Madau (2017). Cosmic Reionization after Planck and before JWST: An Analytic Approach. bibcode 2017ApJ...851...50M

\end{document}
